\section{Impacto social}

\subsection{Una herramienta de primera observación}

Para un pequeño productor, el valor posible de Doctor Maíz no está en reemplazar
un laboratorio ni en emitir una receta. Está en reducir la distancia entre
notar un síntoma y formular una pregunta útil. Una foto puede sugerir si el
patrón se parece más a roya, tizón, daño de gusano cogollero o una deficiencia,
mostrar alternativas y conservar un registro para enseñarlo a un extensionista.
La detección temprana solo ocurre si la captura se hace cuando el síntoma es
visible y si la respuesta conduce a una revisión posterior.

La operación local cambia la viabilidad en zonas con conexión irregular. El
modelo móvil no necesita enviar la fotografía para inferir y WatermelonDB
conserva historial y correcciones hasta que exista conectividad. Esto reduce
consumo de datos y permite trabajar en una parcela sin señal. También concentra
el costo en el teléfono: cámara, memoria y tiempo de inferencia deben ser
aceptables en dispositivos de gama baja, algo que todavía no se midió con una
muestra representativa de equipos.

\subsection{Productores y extensionistas}

Usado individualmente, el sistema puede ayudar a organizar observaciones:
fotografiar varias hojas, comparar confianza, anotar una corrección y decidir
si hace falta asistencia. Usado por extensionistas, puede funcionar como una
ficha común para priorizar visitas y recopilar ejemplos difíciles. El segundo
uso parece más seguro porque una persona con experiencia interpreta el
resultado junto con etapa del cultivo, distribución del daño y condiciones del
suelo.

La cobertura de nueve clases limita ambas funciones. El maíz puede presentar
otras enfermedades, daños por herbicida, sequía, plagas o síntomas combinados.
Una categoría ausente no se convierte en visible por tener un detector OOD.
El usuario necesita una salida clara para «ninguna de las anteriores» y una
ruta de consulta. De lo contrario, la facilidad de tomar una foto puede
producir más confianza de la que permite el dataset.

\subsection{Beneficio posible y evidencia disponible}

El proyecto no midió aumento de rendimiento, ahorro de fertilizante, reducción
de pérdidas ni tiempo de respuesta de asistencia técnica. Atribuirle cualquiera
de esos porcentajes sería inventar impacto. Lo que sí se comprobó es una cadena
técnica: clasificación sobre un corpus amplio, evaluación separada, análisis
por condiciones, inferencia local y mecanismo de corrección. Esa cadena hace
posible una prueba de campo; todavía no demuestra un beneficio agronómico.

Una evaluación de impacto debería comparar prácticas, no solamente modelos.
Por ejemplo, se pueden asignar comunidades o extensionistas a tres condiciones:
procedimiento habitual, acceso a Doctor Maíz y acceso a Doctor Maíz con revisión
técnica. Las medidas relevantes serían tiempo hasta la consulta, concordancia
con diagnóstico agronómico, decisiones evitadas por baja confianza y tasa de
seguimiento. Rendimiento e insumos solo deberían medirse con un diseño que
controle clima, variedad y manejo.

\subsection{Riesgos de exclusión}

El lenguaje, la alfabetización digital y el diseño del teléfono influyen tanto
como el clasificador. Mostrar nombres en inglés, cifras sin explicación o
pantallas cargadas puede excluir a la población a la que se busca apoyar. La
aplicación ya presenta etiquetas y recomendaciones en español, pero no se
realizó una prueba de usabilidad con productores. Tampoco se evaluaron
variantes lingüísticas, accesibilidad visual o instrucciones por audio.

El dataset introduce otra exclusión: proviene de colecciones externas y no
identifica municipios salvadoreños. Un macro-F1 alto puede convivir con fallos
en una variedad local, una cámara económica o una iluminación no representada.
Por eso el despliegue responsable empieza pequeño, registra contexto con
consentimiento y permite detener el uso si una región muestra una brecha
persistente.

\subsection{Condiciones para uso real}

Antes de promover el sistema a escala se necesitan, como mínimo: validación
prospectiva en parcelas; revisión de etiquetas por profesionales; protocolo de
captura sencillo; prueba en teléfonos de distintas gamas; monitoreo por región
y clase; canal de corrección; política de datos y ubicación; y un responsable
que revise incidentes. Bajo esas condiciones, Doctor Maíz puede ampliar el
alcance de la asistencia. Sin ellas, solo traslada al teléfono la incertidumbre
del dataset.
